% Analisis exploratorio de imagenes del proyecto PSIM
% Compilar: latexmk -pdf -outdir=build/reports reports/informe.tex

\documentclass[11pt,a4paper]{article}

\usepackage[utf8]{inputenc}
\usepackage[T1]{fontenc}
\usepackage[spanish]{babel}
\usepackage{lmodern}
\usepackage{geometry}
\geometry{left=2.6cm,right=2.6cm,top=2.8cm,bottom=2.6cm,headheight=1.4cm}
\usepackage{setspace}
\onehalfspacing

\usepackage{hyperref}
\hypersetup{colorlinks=true,linkcolor=blue,citecolor=blue,urlcolor=blue}
\usepackage{graphicx}
\usepackage{fancyhdr}
\usepackage{amsmath,amssymb}
\usepackage{booktabs}
\usepackage{caption}
\usepackage{float}
\usepackage{longtable}
\usepackage{array}

\pagestyle{fancy}
\fancyhf{}
\fancyhead[L]{\includegraphics[height=0.95cm]{Escuela_Rosario_logo.png}}
\fancyfoot[C]{\thepage}
\renewcommand{\headrulewidth}{0pt}

% Las rutas se resuelven desde el directorio de compilacion, que latexmk fija
% en la raiz del proyecto. Declararlas aqui permite compilar con el comando
% publico sin depender de desde donde se invoque.
\makeatletter
\providecommand{\input@path}{}
\edef\input@path{{reports/}{reports/tablas/}{./}}
\makeatother
\graphicspath{{reports/}{figures/fase1/}{figures/fase2/}{figures/fase3/}{./}}

\input{cifras}

\newcommand{\tituloTrabajo}{Analisis exploratorio de imagenes y tabla de caracteristicas}
\newcommand{\subtituloTrabajo}{Estabilidad de los descriptores morfologicos del campo pulmonar}

\begin{document}

\thispagestyle{empty}
\begin{titlepage}
    \centering
    {\LARGE \textbf{\tituloTrabajo}}\\[0.6cm]
    {\large \subtituloTrabajo}\\[2cm]
    {\large \textbf{Juanita Trujillo Narvaez}. Cod: 1000099467}\\
    {\large \textbf{Karol Isabella Martinez Villarreal}. Cod: 1000099682}\\[1cm]
    {\large Ingenieria Biomedica}\\[0.4cm]
    {\large PSIM - 2026II - 80}\\
    Docente: Ing. Pablo Eduardo Caicedo-Rodriguez. Ph.D.
    \vfill
    Fecha de entrega: \today
\end{titlepage}

\tableofcontents
\newpage

\begin{abstract}
Este documento presenta el analisis exploratorio de imagenes del proyecto
semestral. Se audita la procedencia y la calidad de \NImagenesTotal{}
radiografias de torax de dos sitios de adquisicion, se congela una particion
por unidad independiente, se comparan nueve estrategias de segmentacion
clasica contra mascaras de referencia trazadas por personas y se construye
una tabla de \NFilasTabla{} filas con cinco descriptores morfologicos
adimensionales, su valor de referencia y su error relativo. El resultado
principal es negativo y esta sustentado: ninguna estrategia de realce mejora
la linea base sin realce de forma consistente en los dos origenes.
\end{abstract}

\newpage

\section{Problema, objetivos y tarea}

\subsection{Problema}

La delimitacion del campo pulmonar en una radiografia de torax es el paso
previo a cualquier medicion cuantitativa, y es inestable. El contraste entre
el pulmon y las estructuras vecinas depende de la tecnologia del detector, de
los parametros de exposicion y de la presencia de hallazgos que opacifican el
parenquima. Si la mascara obtenida se desplaza segun la estrategia de
procesamiento aplicada, los descriptores morfologicos heredan esa variabilidad
y dejan de ser comparables entre sujetos o entre sitios, aun cuando la
anatomia subyacente no haya cambiado.

\subsection{Decision de modelado}

\begin{table}[H]
\centering
\small
\caption{Ficha de la decision de modelado.}
\begin{tabular}{@{}p{4.0cm}p{10.5cm}@{}}
\toprule
\textbf{Campo} & \textbf{Contenido} \\
\midrule
Objetivo & Determinar que combinacion de realce y umbralizacion produce descriptores con menor error relativo frente a la referencia, de forma consistente entre sitios \\
\addlinespace
Poblacion & \NImagenesTotal{} radiografias PA de Montgomery County (CR) y Shenzhen Hospital (DR), de las cuales \NConMascara{} tienen mascara de referencia \\
\addlinespace
Unidad de analisis & Una imagen procesada bajo una estrategia. La clave es \texttt{sample\_id}, que combina \texttt{image\_id} y \texttt{config\_id} \\
\addlinespace
Unidad independiente & La imagen. Ningun conjunto incluye identificador de paciente; se agrupan solo los dos casos demostrados por referencia cruzada \\
\addlinespace
Entrada disponible & La radiografia PA y sus metadatos tecnicos. La mascara de referencia no existe al procesar una imagen nueva \\
\addlinespace
Salida esperada & Espacial: la mascara del campo pulmonar. De ella se derivan cinco medidas adimensionales \\
\addlinespace
Referencia & Mascaras trazadas por personas. Montgomery bajo supervision de radiologo; Shenzhen por el Instituto Politecnico de Kiev \\
\addlinespace
Metodo previsto & Cadena clasica de realce, umbralizacion y morfologia. Sin etapa supervisada: no se ajusta ningun parametro a partir de etiquetas \\
\addlinespace
Evaluacion & Dice y Jaccard para la mascara; error relativo por descriptor frente a la referencia, estratificado por origen \\
\bottomrule
\end{tabular}
\end{table}

La tarea es de segmentacion. La tabla de caracteristicas acompana al metodo y
no lo sustituye: la salida espacial se conserva como mascara vinculada. No
existe una columna objetivo en sentido supervisado, y la etiqueta de
tuberculosis se registra como variable de control para la comparacion entre
grupos, nunca como objetivo a predecir.

\section{Datos del proyecto}

\subsection{Procedencia y calidad}

\input{inventario}

La caracterizacion de la Fase 1 produjo cinco hallazgos que contradicen o
precisan la documentacion publicada de los conjuntos:

\begin{enumerate}
    \item Las imagenes distribuidas de Montgomery son \texttt{uint8} con ocho
    bits efectivos, no doce como indica la ficha original.
    \item Ninguna de las \NImagenesTotal{} imagenes presenta la escala de
    grises invertida. La inversion que la documentacion atribuye a Shenzhen no
    esta en los archivos distribuidos.
    \item Montgomery mezcla dos orientaciones: 97 verticales y 41
    horizontales. Shenzhen tambien, con 368 y 294.
    \item Shenzhen presenta 525 dimensiones distintas entre 662 imagenes.
    \item Las nueve imagenes mas pequenas de Shenzhen corresponden a los nueve
    pacientes menores de dos anos: el campo de radiografia se ajusta al tamano
    del paciente, y no es un problema de codificacion.
\end{enumerate}

\begin{figure}[H]
\centering
\includegraphics[width=\textwidth]{galeria_entrenamiento.png}
\caption{Galeria de la cohorte de entrenamiento. La regla de muestreo cubre
las dos clases y los dos origenes con dos casos cada uno, los extremos de
resolucion, un paciente pediatrico y un caso marcado para revision en el
inventario. La seleccion dentro de cada criterio es aleatoria con semilla fija.}
\end{figure}

\begin{figure}[H]
\centering
\includegraphics[width=\textwidth]{dimensiones.png}
\caption{Distribucion de dimensiones por origen sobre las \NImagenesTrain{}
imagenes de entrenamiento. Montgomery entrega siempre el mismo tamano de
archivo; Shenzhen lo ajusta al paciente.}
\end{figure}

\subsection{Particion}

\input{particion}

La particion se define sobre unidades independientes y se congela con semilla
fija. Las imagenes que las lecturas clinicas declaran de un mismo paciente
permanecen juntas. La estratificacion es por origen y clase; la edad no entra
en ella porque con cuarenta pacientes menores de dieciocho anos sobre
\NImagenesTotal{} imagenes una tercera variable dejaria celdas de uno o dos
elementos. Su reparto se verifico despues sin forzarlo: la edad mediana es de
34, 35 y 31 anos en entrenamiento, validacion y prueba.

El metodo del proyecto no ajusta parametros a partir de etiquetas, de modo que
no existe el mecanismo por el cual la informacion de prueba podria filtrarse a
un modelo entrenado. La particion cumple aqui dos funciones mas modestas:
acotar el analisis exploratorio al subconjunto de entrenamiento, y dejar
trazable que imagenes se usaron para decidir.

\section{Metodos}

\subsection{Lectura y escalas}

Las radiografias se leen sin alterar sus valores. Las de Shenzhen vienen
codificadas en tres canales RGB identicos, de modo que se toma un solo canal
en lugar de aplicar pesos de luminancia: convertir con aritmetica una imagen
que ya es gris introduce error de redondeo sin aportar nada. Esa identidad se
verifica archivo por archivo y no se asume.

Las intensidades son valores almacenados, no calibrados fisicamente. Ningun
conjunto documenta transformacion de modalidad, y solo Montgomery declara
espaciado de pixel. Por esa razon los descriptores del proyecto son
adimensionales y las unidades espaciales se reportan en pixeles.

\subsection{Region de interes}

Montgomery entrega siempre 4020 por 4892 pixeles con independencia de cuanto
colimo el tecnico. El area no irradiada queda registrada como negro puro y
alcanza el 67 por ciento de la imagen en el caso mas extremo, frente a una
mediana del uno por ciento en Shenzhen. Dividir el area pulmonar por el area
de la imagen mediria entonces la apertura del colimador y no la anatomia del
paciente.

\begin{figure}[H]
\centering
\includegraphics[width=\textwidth]{recorte_verificacion.png}
\caption{Arriba, las seis imagenes de Montgomery con mas area no irradiada,
con el rectangulo envolvente detectado; abajo, el resultado de recortar. El
recorte reduce el fondo pero deja entre un 20 y un 32 por ciento, porque el
area irradiada no es rectangular.}
\end{figure}

\begin{figure}[H]
\centering
\includegraphics[width=0.78\textwidth]{diagnostico_mascara_campo.png}
\caption{Delimitacion del campo irradiado como region. La fila central muestra
el umbral crudo y la inferior la region final. Las zonas oscuras laterales son
las axilas, que atenuan tan poco que se registran al nivel del fondo; no se
rellenan, porque incorporarlas haria que el denominador midiera la apertura
del colimador en lugar del cuerpo del paciente.}
\end{figure}

La region de interes es por tanto la region con senal, obtenida por umbral
sobre el nivel de fondo y seleccion de la componente conexa mayor, sin relleno
de huecos ni cierre morfologico. Sobre ella se calculan el histograma, el
umbral de segmentacion y el denominador de la fraccion de area.

\subsection{Mascaras de referencia}

Los dos conjuntos entregan la referencia en formatos distintos. Montgomery
usa dos archivos booleanos por imagen, uno por pulmon, a la resolucion de la
imagen. Shenzhen usa un archivo \texttt{uint8} con los dos campos, fijo a 512
por 512 pixeles y con valores intermedios propios del antialiasing del borde.

La mascara de Shenzhen se reescala a la imagen por vecino mas cercano, porque
es una etiqueta y no una medida: un valor intermedio entre pulmon y fondo no
tiene interpretacion. Los valores intermedios son el 0,14 por ciento de los
pixeles y binarizar en cualquier umbral entre 1 y 255 cambia el area menos del
0,3 por ciento, de modo que se fija en 128 por reproducibilidad.

El reescalado no crea detalle que la anotacion no tenga. La mascara de
Shenzhen conserva la precision de 512 por 512 aunque se represente sobre una
rejilla mayor, y esa limitacion se declara sin compensarla.

\subsection{Estrategias comparadas}

Se comparan nueve combinaciones de tres realces y tres umbralizaciones, con la
morfologia fija en todas. Dejarla constante es deliberado: si variara junto
con lo demas, una diferencia en el resultado no podria atribuirse a una causa
concreta. La secuencia morfologica es cierre, relleno de huecos, apertura y
seleccion de las dos componentes mayores.

\begin{table}[H]
\centering
\small
\caption{Parametros registrados de cada familia de tecnicas.}
\begin{tabular}{@{}p{3.6cm}p{4.2cm}p{6.4cm}@{}}
\toprule
\textbf{Familia} & \textbf{Decision} & \textbf{Parametros} \\
\midrule
Muestreo y color & Aplicar & Canal unico verificado; lado de trabajo 512 px; sin espaciado fisico \\
\addlinespace
Intensidad e histogramas & Comparar & Sin realce, ecualizacion por CDF y CLAHE con limite 0,02 y rejilla de 8 \\
\addlinespace
Vecindarios y gradiente & No pertinente & El borde pulmonar se delimita por umbral; el gradiente no entra en el flujo final \\
\addlinespace
Umbralizacion y morfologia & Comparar & Umbral fijo al percentil 35, Otsu y adaptativo con bloque de 0,25 del lado; apertura 0,012 y cierre 0,020 \\
\addlinespace
Fourier 2D & No pertinente & No se observaron patrones periodicos que justifiquen filtrado en frecuencia \\
\addlinespace
Wavelets 2D & Aplicar & Familia db2, nivel 1, energia de los tres detalles normalizada por area \\
\bottomrule
\end{tabular}
\end{table}

Las tres estrategias sin realce son la linea base contra la que se mide el
criterio de exito, y se conservan en todos los resultados.

\subsection{Descriptores}

\begin{table}[H]
\centering
\small
\caption{Matriz de objetivo, propiedad observable, descriptor y uso previsto.}
\begin{tabular}{@{}p{3.4cm}p{3.6cm}p{3.4cm}p{3.6cm}@{}}
\toprule
\textbf{Objetivo} & \textbf{Propiedad observable} & \textbf{Descriptor} & \textbf{Uso previsto} \\
\midrule
Cuantificar la forma del campo & Regularidad del contorno & Compacidad & Estabilidad frente a la estrategia \\
\addlinespace
Cuantificar la proporcion & Extension vertical frente a horizontal & Relacion de aspecto & Estabilidad y control de edad \\
\addlinespace
Cuantificar el volumen aireado & Area pulmonar sobre el cuerpo & Fraccion de area & Estabilidad y comparacion entre grupos \\
\addlinespace
Detectar afectacion unilateral & Diferencia entre campos & Simetria izq-der & Comparacion entre grupos \\
\addlinespace
Caracterizar el parenquima & Textura a escala fina & Energia wavelet & Comparacion entre origenes \\
\bottomrule
\end{tabular}
\end{table}

Cada descriptor se calcula igual sobre la mascara obtenida y sobre la de
referencia, de modo que el error relativo entre ambos valores sea
interpretable. Ninguna funcion de extraccion recibe la etiqueta de
tuberculosis.

\section{Resultados}

\subsection{Variabilidad entre sitios}

\begin{figure}[H]
\centering
\includegraphics[width=\textwidth]{intensidad.png}
\caption{Histogramas por imagen y resumen por unidad de analisis sobre las
\NImagenesTrain{} imagenes de entrenamiento. El panel central muestra que los
dos origenes forman nubes casi disjuntas en mediana e intervalo
intercuartilico.}
\end{figure}

La mediana de intensidad es de 111,0 en Montgomery y 173,5 en Shenzhen, con
intervalos intercuartilicos de 151,5 y 90,0. Dentro de cada origen, en cambio,
las clases apenas difieren: en Montgomery la mediana es 110,0 para los casos
sin hallazgos y 111,0 para los casos con manifestaciones.

El efecto del sitio de adquisicion supera por tanto al de la clase. Un
descriptor de intensidad que separase las clases estaria separando hospitales,
y esa es la razon por la que toda metrica se reporta estratificada por origen.

\subsection{Concordancia con la referencia}

\input{dice}

\begin{figure}[H]
\centering
\includegraphics[width=\textwidth]{dice_por_estrategia.png}
\caption{Concordancia por estrategia. En verde la linea base sin realce. El
panel derecho contrasta el valor mediano en cada origen.}
\end{figure}

La mejor estrategia es \MejorEstrategia{}, con Dice mediano de \MejorDice{}, y
es ademas la mas consistente entre sitios: \MejorDiceMC{} en Montgomery y
\MejorDiceSH{} en Shenzhen. No aplica ningun realce.

\subsection{Desplazamiento de los descriptores}

\input{errores}

\begin{figure}[H]
\centering
\includegraphics[width=\textwidth]{error_descriptores.png}
\caption{Error relativo mediano de cada descriptor frente a su valor de
referencia, por estrategia. En verde la linea base.}
\end{figure}

La relacion de aspecto es el descriptor mas estable, con error relativo
mediano entre 0,066 y 0,158 segun la estrategia. La simetria izquierda-derecha
es el menos estable, entre 0,647 y 0,839. La compacidad se situa entre 0,410 y
0,593.

Parte de la inestabilidad de la simetria tiene una causa identificable: una
segmentacion que pierde un campo pulmonar produce asimetria maxima, de modo
que el descriptor mezcla variacion anatomica con fallo del metodo.

\subsection{Criterio de exito}

El criterio fijado antes de procesar los datos exige al menos una combinacion
cuyo error relativo mediano mejore la linea base en tres de los cinco
descriptores, de forma consistente en ambos origenes.

\begin{table}[H]
\centering
\small
\caption{Evaluacion del criterio de exito. Cada estrategia con realce se
compara contra la linea base que comparte umbralizacion.}
\begin{tabular}{@{}llccl@{}}
\toprule
\textbf{Estrategia} & \textbf{Linea base} & \textbf{Montgomery} & \textbf{Shenzhen} & \textbf{Resultado} \\
\midrule
clahe\_adap & base\_adap & 0/5 & 2/5 & no cumple \\
clahe\_fijo & base\_fijo & 2/5 & 4/5 & no cumple \\
clahe\_otsu & base\_otsu & 1/5 & 2/5 & no cumple \\
eq\_adap & base\_adap & 3/5 & 1/5 & no cumple \\
eq\_fijo & base\_fijo & 0/5 & 2/5 & no cumple \\
eq\_otsu & base\_otsu & 2/5 & 3/5 & no cumple \\
\bottomrule
\end{tabular}
\end{table}

Ninguna de las seis combinaciones con realce cumple el criterio. Las dos que
se acercan fallan precisamente por inconsistencia entre sitios: clahe\_fijo
mejora cuatro de cinco descriptores en Shenzhen pero solo dos en Montgomery, y
eq\_adap lo contrario.

Ese resultado es negativo y esta sustentado. La pregunta de investigacion
indaga en que medida las estrategias modifican la estabilidad de los
descriptores, no presupone que la mejoren, y la tabla de riesgos del proyecto
contemplaba este escenario antes de procesar los datos. Si se hubiera usado un
promedio agregado en lugar de exigir consistencia en ambos origenes,
clahe\_fijo habria parecido una mejora.

\section{Decisiones para el modelado}

\input{decisiones}

\section{Reproducibilidad y limites}

\subsection{Tabla de caracteristicas}

La tabla tiene \NFilasTabla{} filas y \NColumnasTabla{} columnas: una fila por
cada una de las \NConMascara{} imagenes con referencia en entrenamiento,
procesada bajo cada una de las \NEstrategias{} estrategias. La clave es
\texttt{sample\_id}.

Los cinco descriptores calculados sobre la mascara obtenida son los unicos
predictores declarados en \texttt{config/model\_input.json}. Las columnas con
prefijo \texttt{ref\_} y \texttt{err\_}, junto con el Dice y el Jaccard,
quedan explicitamente excluidas: derivan de la mascara de referencia, que no
existe al procesar una imagen nueva. Usarlas como predictores seria circular.

\input{auditoria}

\subsection{Reconstruccion}

Desde una copia sin entorno virtual, la secuencia publica es:

\begin{verbatim}
uv sync --locked
uv run ruff check .
uv run ruff format --check .
uv run pytest
uv run python scripts/reproduce.py
latexmk -pdf -outdir=build/reports reports/informe.tex
latexmk -pdf -outdir=build/slides slides/presentacion.tex
\end{verbatim}

Ningun numero de este documento se transcribe a mano: las tablas se generan
desde los archivos de resultados y las cifras citadas en el texto proceden de
macros producidas por el mismo guion. Los datos originales no se versionan por
sus condiciones de uso; \texttt{data/raw/README.md} documenta como obtenerlos
y \texttt{scripts/check\_dataset.py} verifica que la copia local corresponde al
conjunto descrito.

\subsection{Limites de interpretacion}

\begin{enumerate}
    \item Las mascaras de referencia proceden de dos trazados con convenciones
    distintas. Las de Montgomery excluyen corazon y grandes vasos; las de
    Shenzhen los delimitan. Es un sesgo sistematico entre origenes que no
    puede corregirse.
    \item Ningun conjunto incluye identificador de paciente. Se agrupan los
    dos casos demostrados por referencia cruzada, y 90 imagenes de Montgomery
    comparten sexo y edad con alguna otra sin que eso pruebe identidad.
    \item El resultado sobre el realce es valido para los parametros
    declarados. Un CLAHE con otro limite de recorte podria comportarse de otro
    modo; ese espacio no se exploro.
    \item Las 96 imagenes de Shenzhen sin mascara quedan fuera de toda
    comparacion contra referencia. Participan en la caracterizacion pero no en
    la evaluacion.
    \item Con dos sitios no se puede distinguir si la diferencia observada
    proviene del detector, del protocolo o de la poblacion.
\end{enumerate}

\subsection{Contribuciones}

Juanita Trujillo mantuvo la infraestructura reproducible, el entorno y los
guiones de verificacion e inventario. Karol Martinez mantuvo el documento y la
coherencia entre los resultados y la narrativa. Cada bloque tecnico tuvo una
integrante que lo lidero y otra que lo reviso ejecutando el codigo.

\nocite{*}
\bibliographystyle{plain}
\bibliography{referencias}

\end{document}